\section{Resultados Esperados}

De acuerdo con las brechas identificadas en el diagnóstico inicial de madurez digital y los objetivos estratégicos definidos para el proyecto, la implementación del sistema Odoo se orientó a lograr una transformación estructural en los procesos operativos y comerciales de la agencia. Se establecieron los siguientes resultados esperados:

\begin{itemize}
    \item \textbf{Centralización y Seguridad de la Información:} Se proyectó la eliminación definitiva de los "silos de información" mediante una base de datos centralizada en la nube que integre perfiles 360° de pasajeros y proveedores. Esto permite erradicar la dispersión de datos en Excel, WhatsApp y cuadernos físicos, garantizando la trazabilidad integral de cada servicio.
    \item \textbf{Optimización del Tiempo de Respuesta Comercial:} El resultado más crítico es la reducción del tiempo promedio de elaboración de cotizaciones complejas, pasando de 45-60 minutos a una respuesta automatizada en menos de 15 minutos. Esto busca dotar a la agencia de una agilidad competitiva frente a operadores digitalizados y OTAs.
    \item \textbf{Estandarización de Paquetes Turísticos:} Se esperaba digitalizar el catálogo de servicios (Cusco, Arequipa, Puno), permitiendo la generación instantánea de itinerarios y presupuestos profesionales en PDF. La automatización del flujo comercial busca eliminar la necesidad de copiar y pegar información manualmente entre documentos.
    \item \textbf{Reducción Drástica del Riesgo Operativo:} A través de la conversión directa de cotizaciones a órdenes de venta, se proyectó reducir la tasa de error operativo (en nombres, fechas y servicios) del 5\% a menos del 1\%. Esto impacta directamente en la rentabilidad al evitar reprocesos y mejorar la calidad percibida por el viajero.
    \item \textbf{Gestión Sistemática de la Calidad:} Implementar un sistema de retroalimentación post-viaje mediante el módulo de Encuestas, permitiendo medir el nivel de satisfacción (NPS) de manera automatizada. El objetivo es capturar el "boca a boca" digital para fortalecer la reputación de la marca en el sitio web.
    \item \textbf{Disponibilidad de Datos para Decisiones:} Finalmente, se esperaba que la gerencia contara con tableros de control (dashboards) en tiempo real para visualizar la rentabilidad por paquete y el desempeño de ventas, facilitando una dirección estratégica basada en datos reales.
\end{itemize}
